% Preámbulo común de los diagramas del Objetivo 3 (paleta de main.tex).
\usepackage[T1]{fontenc}
\usepackage[utf8]{inputenc}
\usepackage{lmodern}
\usepackage[spanish]{babel}
\usetikzlibrary{babel,arrows.meta,positioning,fit,backgrounds,shapes.multipart,calc}
\definecolor{dPersonaF}{HTML}{FFE6CC}\definecolor{dPersonaB}{HTML}{D79B00}
\definecolor{dSistemaF}{HTML}{DAE8FC}\definecolor{dSistemaB}{HTML}{6C8EBF}
\definecolor{dAgenteF}{HTML}{D5E8D4}\definecolor{dAgenteB}{HTML}{82B366}
\definecolor{dDatosF}{HTML}{E1D5E7}\definecolor{dDatosB}{HTML}{9673A6}
\definecolor{dClinicoF}{HTML}{FFF2CC}\definecolor{dClinicoB}{HTML}{D6B656}
\definecolor{dExternoF}{HTML}{F5F5F5}\definecolor{dExternoB}{HTML}{666666}
% Entidad: encabezado coloreado y columnas clave. PK subrayada, FK en cursiva.
\tikzset{
  ent/.style 2 args={rectangle split, rectangle split parts=2, draw=#2, line width=0.7pt,
    rectangle split part fill={#1,white}, rounded corners=1.5pt, font=\footnotesize,
    inner sep=3pt, minimum width=2.9cm},
  ext/.style={draw=dExternoB, dashed, fill=dExternoF, rounded corners=1.5pt, font=\footnotesize\ttfamily,
    inner sep=4pt, minimum width=2.4cm},
  rel/.style={draw=black!65, line width=0.6pt},
  card/.style={font=\scriptsize, inner sep=1.5pt, fill=white},
}
\newcommand{\pk}[1]{\underline{#1}}
\newcommand{\fk}[1]{\textit{#1}}
\newcommand{\cols}[1]{\begin{tabular}{@{}l@{}}#1\end{tabular}}
\newcommand{\tname}[1]{\textbf{\ttfamily #1}}
% Relación 1:N con cardinalidades junto a cada extremo.
\newcommand{\unoN}[3][]{\draw[rel,#1] (#2) -- node[card, pos=0.12]{1} node[card, pos=0.88]{N} (#3);}
\newcommand{\unoUno}[3][]{\draw[rel,#1] (#2) -- node[card, pos=0.12]{1} node[card, pos=0.88]{1} (#3);}
\newcommand{\leyenda}{\node[font=\scriptsize, text=black!70, anchor=north west] at (current bounding box.south west)
  {\underline{subrayado}: clave primaria \quad \textit{cursiva}: clave foránea \quad recuadro gris: tabla de otro grupo};}
